\begin{definition}[Archimedean Invariants of a Number Field] \label{definition:archimedean_invariants_of_a_number_field}
Let $K$ be a \CrefAndHyperrefIfExist{definition:number_field}{number field} of degree $n = [K : \mathbb{Q}]$. Let $\sigma_1, \ldots, \sigma_n : K \hookrightarrow \mathbb{C}$ be the distinct embeddings of $K$ into $\mathbb{C}$. Since $K$ is a subfield of $\mathbb{C}$ and complex conjugation is an automorphism of $\mathbb{C}$, it acts on the set of these \CrefAndHyperrefIfExist{definition:field_homomorphism}{embeddings}.

The number of embeddings $\sigma_i$ whose image lies in $\mathbb{R}$ is denoted by \hl{$r_1$} and called the \hldef{number of real embeddings} (or the \hldef{real signature}) of $K$. The number of pairs of distinct complex conjugate embeddings among $\{\sigma_1, \ldots, \sigma_n\}$ is denoted by \hl{$r_2$} and called the \hldef{number of complex embedding pairs} (or the \hldef{complex signature}) of $K$.

The pair $(r_1, r_2)$ is called the \hldef{signature of $K$}. These invariants satisfy
$$n = r_1 + 2r_2.$$
\end{definition}
